\section{Evaluation}
\label{sec:eval}

\subsection{Models}

We evaluate two instruction-tuned LLMs: \qwen{} (Qwen3-Coder-30B-A3B-Instruct,
served locally via LM Studio) and \nemotron{} (nvidia/Nemotron-3-Ultra-550B-A55B,
accessed via the NVIDIA Inference API with \texttt{enable\_thinking: false}).
Both support OpenAI-compatible function calling.

\subsection{RQ1: Does Graph Context Improve Accuracy?}

Table~\ref{tab:accuracy} reports exact and top-3 accuracy per condition.
\emph{Exact} counts runs where \texttt{answer\_function} matches the ground
truth; \emph{top-3} counts questions where the ground truth appears in
\texttt{top3} or equals \texttt{answer\_function}.
Q2's ground truth (\texttt{sqlite3Select}) was correct in all three \qwen{}
\textsc{NOCTX} repetitions, suggesting memorisation; we flag and exclude Q2
from \qwen{} counts in the primary analysis ($n{=}5$), and also report $n{=}6$
to allow direct comparison with \nemotron{} ($n{=}6$, no questions flagged).

\begin{table}[h]
\caption{Accuracy per condition.
  Qwen primary: $n=5$ (Q2 excluded). Qwen full: $n=6$ (Q2 included).
  Nemotron: $n=6$.
  Dashes: ablation not run on this model.
  NOCTX top-3 counts questions, not runs (3~reps each).}
\label{tab:accuracy}
\small
\begin{tabular}{lrr|rr|rr}
\toprule
 & \multicolumn{2}{c|}{\textbf{Qwen ($n{=}5$)}}
 & \multicolumn{2}{c|}{\textbf{Qwen ($n{=}6$)}}
 & \multicolumn{2}{c}{\textbf{Nemotron ($n{=}6$)}} \\
\cmidrule(lr){2-3}\cmidrule(lr){4-5}\cmidrule(lr){6-7}
\textbf{Condition} & \textbf{Ex.} & \textbf{T3}
                   & \textbf{Ex.} & \textbf{T3}
                   & \textbf{Ex.} & \textbf{T3} \\
\midrule
\textsc{NOCTX}            & 0 & 0 & 1 & 1 & 1 & 2 \\
\textsc{SYN-TOOL}          & 1 & 2 & 1 & 3 & 1 & 4 \\
\textsc{CT-TOOL}           & 0 & 1 & 0 & 2 & 1 & 3 \\
\textsc{SYN-TEXT}          & 1 & 1 & 1 & 1 & 1 & 3 \\
\textsc{CT-TEXT}           & 1 & 1 & 2 & 2 & 1 & 2 \\
\midrule
\textsc{CT-ct\_calls}      & 0 & 1 & --- & --- & --- & --- \\
\textsc{CT-ct\_types}      & 0 & 1 & --- & --- & --- & --- \\
\textsc{CT-src-only}       & 0 & 1 & --- & --- & --- & --- \\
\textsc{SYN-src-only}      & 0 & 1 & --- & --- & --- & --- \\
\bottomrule
\end{tabular}
\end{table}

Overall accuracy is low across all conditions: no condition exceeds one exact
answer out of five (Qwen $n{=}5$) or one out of six (Nemotron) in the primary
analysis.
Both models score near zero in the no-context baseline: 1 of 15 \qwen{}
\textsc{NOCTX} runs exact (Q2, flagged); 1 of 18 \nemotron{} \textsc{NOCTX}
runs exact (Q1).
Providing graph context does not reliably improve exact accuracy; on some
questions it adds noise by giving the model structurally incomplete information
to reason over.

\subsection{RQ2: Tool vs.\ Text Slice}

For \qwen{} ($n{=}5$), \textsc{SYN-TEXT} and \textsc{CT-TEXT} each score one
exact answer, matching \textsc{SYN-TOOL} and exceeding \textsc{CT-TOOL} (zero
exact).
For \nemotron{}, all four graph conditions score one exact answer;
\textsc{SYN-TOOL} achieves the highest top-3 count (4 of 6).
\nemotron{} solved Q5 (\texttt{sqlite3\_str\_vappendf}) in both tool conditions
while \qwen{} succeeded only in the text conditions, so the text-over-tool
pattern on Q5 holds for \qwen{} but not across both models.

Note that TOOL and TEXT conditions differ in more than delivery mode: in the
TOOL conditions the agent can call \texttt{get\_function\_source} to read raw
source; in the TEXT conditions the slice is the only source representation and
raw source is not separately available.
This confound is discussed further in Section~\ref{sec:threats}.

\paragraph{Token cost (Table~\ref{tab:cost}).}
Text slices consume 3,500--5,500 mean prompt tokens.
Tool conditions cost roughly 10--12$\times$ more for \qwen{} and 16--23$\times$
more for \nemotron{} (Nemotron SYN-TOOL: 81,796 tokens; CT-TOOL: 86,373).
Every tool-condition run hit the 15-turn cap and ended with a forced final
answer, so the token budget was fully spent in every case.

\begin{table}[h]
\caption{Mean prompt-token cost and latency by condition.
  NOCTX/SYN-TOOL/CT-TOOL/SYN-TEXT/CT-TEXT rows report Nemotron figures;
  ablation rows report \qwen{} figures.}
\label{tab:cost}
\small
\begin{tabular}{llrr}
\toprule
\textbf{Condition} & \textbf{Model} & \textbf{Mean prompt tokens} & \textbf{Mean latency (s)} \\
\midrule
\textsc{NOCTX}             & Nemotron &         342 &  11 \\
\textsc{SYN-TOOL}           & Nemotron &      81,796 &  54 \\
\textsc{CT-TOOL}            & Nemotron &      86,373 &  48 \\
\textsc{SYN-TEXT}           & Nemotron &       3,534 &  19 \\
\textsc{CT-TEXT}            & Nemotron &       5,480 &   7 \\
\textsc{CT-src-only}        & Qwen     & $\approx$209,000 & 55 \\
\textsc{SYN-src-only}       & Qwen     & $\approx$209,000 & 49 \\
\bottomrule
\end{tabular}
\end{table}

\subsection{Error Analysis}

To understand failure modes, every \qwen{} answer across all eight graph
conditions (48 total answer slots across 6 questions; \textsc{NOCTX} excluded)
was classified.
Five slots contain correct answers; the remaining 43 wrong-answer slots fall
into one of four buckets, presented here in priority order (a slot is assigned
to the first matching bucket).

\paragraph{Callback gap (15/48 slots, 35\,\% of wrong answers).}
The ground-truth function is unreachable from the entry in either graph because
the call crosses a function pointer.
This is the dominant failure for Q1 (7 of 8 conditions wrong) and the exclusive
failure for Q4 (all 8 conditions wrong).
Neither graph records function-pointer edges, so a pure graph traversal along
static edges cannot bridge the gap; the static graph slice also lacks the ground-truth
function for both questions. As shown in Section~\ref{sec:turn_horizon}, bridging this
gap requires the agent to combine symbol search and source reading across an extended
turn horizon.

\paragraph{Wrong other (25/48 slots, 58\,\% of wrong answers).}
The ground-truth function is a graph node, but the model named a different one.
For Q3 (\texttt{flattenSubquery}), the graph contains the ground truth in both
graphs yet the model named \texttt{codeDistinct} across six of eight conditions.
For Q5 (\texttt{sqlite3\_str\_vappendf}) on \qwen, every tool condition stops
at \texttt{printfTempBuf} (which \texttt{sqlite3\_str\_vappendf} calls);
the text slice surfaces the target at shallow BFS depth.

\paragraph{Default to entry (3/48 slots, 7\,\% of wrong answers).}
The model named the entry function itself (\texttt{sqlite3Select} or
\texttt{sqlite3\_step}).
Note that \texttt{sqlite3Select} answers for Q3 and Q6 are classified here
rather than as ``wrong other'' because \texttt{sqlite3Select} is the entry
function for those questions.

\paragraph{Not in graph (0/48 slots).}
The ground-truth function is a node in both graphs for every question that is
not a callback-gap question.

\subsection{RQ4: Path Correctness}

Every cited edge (\texttt{src}, \texttt{dst}, \texttt{kind}) and every
consecutive evidence-path pair from all 36 \qwen{} tool-condition runs were
checked against the graph database.
An edge is recorded as \emph{not supported by the graph} when neither direction
appears; because neither graph models function-pointer calls, an absent edge
may reflect a genuine structural gap rather than model fabrication, so this
rate is an upper bound on hallucination.

Of 133 cited edges, 106 (79\,\%) are present exactly as stated and 4 (3\,\%)
have direction reversed.
The not-in-graph rate was 2.4$\times$ lower in conditions with caller/callee
tools (12\,\%, 10 of 86 edges) than in source-only conditions (28\,\%, 13 of 47
edges).
The pattern is similar for evidence-path pairs (13\,\% vs.\ 34\,\%).
Source-only results suggest the model cites edges it cannot verify when no
edge-querying tool is available.

One further observation: the one \qwen{} run that correctly identified Q1
(\textsc{SYN-TOOL}) still contains one cited edge not supported by the graph,
necessarily spanning the function-pointer gap; a correct final answer does not
guarantee a fully supported reasoning chain.

\subsection{RQ3: Ablation}

No ablation condition scored above zero exact answers on \qwen's five unflagged
questions, and every ablation run hit the 15-turn cap.
Adding callers, callees, type hierarchy, and taint paths to the source-reading
baseline did not improve accuracy on any question.
The source-only control costs about 61$\times$ more tokens than \textsc{SYN-TEXT}
and 40$\times$ more than \textsc{CT-TEXT}, with no accuracy gain.
When source-reading is the only retrieval tool, the model paged through raw
function source across all fifteen turns rather than navigating with callers and
callees; every source-only run ended with a forced final answer.

\subsection{Turn Horizon and Guidance Sensitivity Analysis}
\label{sec:turn_horizon}

In the baseline protocol, all 48 tool runs hit the 15-turn cap and required a
forced final answer.
To understand the underlying causes, we conducted an exploratory sensitivity
analysis examining turn awareness, stopping cues, and search horizon on \qwen.
Per our pre-registered decision log (D17), all modifications are evaluated on a
separate run track and reported side-by-side; they do not replace the headline
baseline figures.

\paragraph{Turn awareness and early stopping (15 turns).}
We updated the tool prompt to provide a per-turn countdown (\emph{``Tool calls left:~$N$''})
and a convergence nudge at two-thirds budget (\emph{``Start converging. If you have a candidate, answer now''}).
Under this 15-turn configuration, 4 of 12 runs concluded with \texttt{stop=final} on their
own without invoking the forced final prompt.

\paragraph{Extended search horizon with 3 repetitions (30 turns, $T=0.7$).}
To evaluate whether performance gains were stable or stochastic, we ran 3 sampled repetitions
per question at temperature~0.7 with a 30-turn cap and the refined guidance prompt (36 runs total).
The results confirm systematic improvements across both graph back-ends:
\begin{itemize}
  \item \textbf{Autonomous termination}: Rather than hitting the turn cap indiscriminately,
    \textbf{26 of 36 runs (72\,\%)} terminated autonomously with \texttt{stop=final},
    frequently concluding between turns 22 and 26.
  \item \textbf{Callback-gap targets reached}: On \textsc{CT-TOOL}, \texttt{resolveOrderByTermToExprList} (Q1)
    was solved in \textbf{2 of 3 repetitions} (67\,\% consistency).
    Crucially, \texttt{withExpand} (Q4)---which was never solved by any condition or model in the
    15-turn baseline---was identified on \textbf{both graphs} (1 of 3 runs on \textsc{SYN-TOOL},
    and 1 of 3 runs on \textsc{CT-TOOL}).
  \item \textbf{Overall accuracy}: \textsc{SYN-TOOL} solved 3 distinct questions across runs (Q2, Q4, and
    Q5 in 2 of 3 runs); \textsc{CT-TOOL} solved 2 distinct questions (Q1 and Q4), with 100\,\% top-3 coverage
    on Q2 and 67\,\% top-3 coverage on Q1 and Q5.
\end{itemize}

\paragraph{Methodological takeaways.}
Tuning guidance prompts on the evaluation benchmark carries an inherent risk of
optimistic bias, and the additional navigation depth incurs substantial token cost
(averaging 250{,}000--400{,}000 prompt tokens per query).
Nonetheless, the 3-repetition experiment provides strong evidence that: (i)~the 15-turn cap was
artificially truncating necessary search depth on multi-hop and callback-separated functions,
(ii)~turn countdown visibility effectively curtails unbounded exploration, and (iii)~static
function-pointer gaps do not represent an insurmountable ceiling when agents can inspect source
and search symbol definitions.
